\documentclass[aspectratio=169,10pt]{beamer}

% =========================================================
% PAQUETES
% =========================================================
\usepackage[spanish,es-noshorthands]{babel}
\usepackage[utf8]{inputenc}
\usepackage[T1]{fontenc}
\usepackage{lmodern}
\usepackage{amsmath,amssymb,mathtools}
\usepackage{array}
\usepackage{booktabs}
\usepackage{multicol}
\usepackage{tikz}
\usepackage{enumitem}
\usepackage{xcolor}

% =========================================================
% TEMA
% =========================================================
\usetheme{Madrid}
\usecolortheme{default}
\setbeamertemplate{navigation symbols}{}

% =========================================================
% COLORES
% =========================================================
\definecolor{USTBlue}{RGB}{0,70,127}
\definecolor{USTBlueDark}{RGB}{0,45,85}
\definecolor{USTTeal}{RGB}{0,153,117}
\definecolor{USTSoft}{RGB}{245,247,250}
\definecolor{USTGray}{RGB}{90,90,90}
\definecolor{USTText}{RGB}{35,40,45}
\definecolor{USTPurple}{RGB}{109,76,139}

\setbeamercolor{title}{fg=white,bg=USTBlue}
\setbeamercolor{frametitle}{fg=white,bg=USTBlue}
\setbeamercolor{structure}{fg=USTBlue}
\setbeamercolor{normal text}{fg=USTText,bg=white}
\setbeamercolor{block title}{fg=white,bg=USTBlueDark}
\setbeamercolor{block body}{fg=black,bg=USTSoft}
\setbeamercolor{block title alerted}{fg=white,bg=USTPurple}
\setbeamercolor{block body alerted}{fg=black,bg=USTSoft}
\setbeamercolor{block title example}{fg=white,bg=USTTeal}
\setbeamercolor{block body example}{fg=black,bg=USTSoft}

% =========================================================
% FOOTLINE
% =========================================================
\setbeamertemplate{footline}{
	\leavevmode
	\hbox{
		\begin{beamercolorbox}[wd=.78\paperwidth,ht=2.6ex,dp=1.2ex,leftskip=0.4cm]{author in head/foot}
			\color{white}\bfseries Guion de Video -- Evaluación Unidad 1 -- Pensamiento Matemático
		\end{beamercolorbox}
		\begin{beamercolorbox}[wd=.22\paperwidth,ht=2.6ex,dp=1.2ex,rightskip=0.35cm plus1fil]{date in head/foot}
			\color{white}\bfseries \insertframenumber/\inserttotalframenumber
		\end{beamercolorbox}
	}
}
\setbeamercolor{author in head/foot}{bg=USTBlueDark,fg=white}
\setbeamercolor{date in head/foot}{bg=USTBlue,fg=white}

% =========================================================
% ENTORNOS
% =========================================================
\newcommand{\IdeaClave}[1]{
	\begin{alertblock}{Idea clave}
		#1
	\end{alertblock}
}

\newcommand{\Guion}[1]{
	\begin{exampleblock}{Guion sugerido para narrar}
		#1
	\end{exampleblock}
}

\newcommand{\Resultado}[1]{
	\begin{block}{Resultado}
		#1
	\end{block}
}

% =========================================================
% DATOS
% =========================================================
\title{Guion Completo en Video\\Evaluación Unidad 1}
\subtitle{Pensamiento Matemático -- Psicología}
\author{Luis Tulio González Alcaino}
\institute{Universidad Santo Tomás \\ Departamento de Ciencias Básicas}
\date{}

% =========================================================
\begin{document}
	% =========================================================
	
	% ---------------------------------------------------------
	\begin{frame}
		\titlepage
	\end{frame}
	
	% ---------------------------------------------------------
	\begin{frame}{Ruta del video}
		\tableofcontents
	\end{frame}
	
	% =========================================================
	\section{Estructura general del video}
	% =========================================================
	
	% ---------------------------------------------------------
	\begin{frame}{Estructura sugerida del video}
		\begin{block}{Duración recomendada}
			12 a 18 minutos.
		\end{block}
		
		\begin{enumerate}[label=\arabic*.]
			\item Presentación inicial.
			\item Problema 1.
			\item Problema 2.
			\item Problema 3.
			\item Problema 4.
			\item Cierre final.
		\end{enumerate}
		
		\IdeaClave{
			En el video se debe mostrar no solo la respuesta final, sino también el procedimiento, la justificación lógica y la interpretación matemática.
		}
	\end{frame}
	
	% ---------------------------------------------------------
	\begin{frame}{Recomendaciones para la grabación}
		\begin{itemize}
			\item Hablar de forma pausada y clara.
			\item Escribir cada paso en hoja, pizarra o diapositiva.
			\item Nombrar la propiedad lógica o la operación de conjuntos que se está usando.
			\item Evitar dar solo el resultado.
			\item Justificar cada conclusión.
		\end{itemize}
		
		\Guion{
			``Hola. En este video resolveré paso a paso la Evaluación de la Unidad 1 de Pensamiento Matemático, contextualizada al área de la Psicología. Desarrollaré cada problema mostrando el procedimiento, la justificación lógica y la interpretación final de los resultados.''}
	\end{frame}
	
	% =========================================================
	\section{Problema 1}
	% =========================================================
	
	% ---------------------------------------------------------
	\begin{frame}{Problema 1: Enunciado}
		Dadas las proposiciones:
		
		\begin{align*}
			p &: \text{La memoria a corto plazo tiene capacidad ilimitada.} \\
			q &: \text{La memoria sensorial tiene una duración de varios minutos.} \\
			r &: \text{La memoria a largo plazo es de duración prolongada.} \\
			s &: \text{El individuo utiliza la información para tomar decisiones.}
		\end{align*}
		
		\begin{enumerate}[label=\alph*)]
			\item Traduzca a lenguaje usual:
			\[
			(p \to q)\land(\sim r \to s)
			\]
			\item Determine el valor de verdad de:
			\[
			[(p\land \sim q)\lor r]\to s
			\]
		\end{enumerate}
	\end{frame}
	
	% ---------------------------------------------------------
	\begin{frame}{Problema 1(a): Traducción a lenguaje usual}
		\Guion{
			``Primero analizo la proposición compuesta \((p\to q)\land(\sim r\to s)\).\\[0.2cm]
			La primera parte, \(p\to q\), se lee: si \(p\), entonces \(q\).\\
			La segunda parte, \(\sim r\to s\), se lee: si no \(r\), entonces \(s\).\\
			Como ambas están unidas por la conjunción \(\land\), la traducción final se expresa usando la palabra \textit{y}.''}
		
		\Resultado{
			\[
			(p\to q)\land(\sim r\to s)
			\]
			se traduce como:
			\vspace{0.2cm}
			
			\emph{``Si la memoria a corto plazo tiene capacidad ilimitada, entonces la memoria sensorial tiene una duración de varios minutos, y si la memoria a largo plazo no es de duración prolongada, entonces el individuo utiliza la información para tomar decisiones.''}
		}
	\end{frame}
	
	% ---------------------------------------------------------
	\begin{frame}{Problema 1(b): Análisis lógico}
		\Guion{
			``Ahora debo determinar el valor de verdad de la proposición
			\[
			[(p\land \sim q)\lor r]\to s.
			\]
			Para ello asigno a cada proposición simple un valor de verdad razonable según el contexto psicológico.''}
		
		\begin{block}{Valores de verdad}
			\[
			p=F,\qquad q=F,\qquad r=V,\qquad s=V
			\]
		\end{block}
		
		\IdeaClave{
			La memoria a corto plazo no tiene capacidad ilimitada, la memoria sensorial no dura varios minutos, la memoria a largo plazo sí es prolongada y el individuo sí usa información para tomar decisiones.
		}
	\end{frame}
	
	% ---------------------------------------------------------
	\begin{frame}{Problema 1(b): Desarrollo paso a paso}
		Sustituyendo en
		\[
		[(p\land \sim q)\lor r]\to s
		\]
		obtenemos:
		\[
		[(F\land \sim F)\lor V]\to V
		\]
		
		Como
		\[
		\sim F = V,
		\]
		queda:
		\[
		[(F\land V)\lor V]\to V
		\]
		
		Luego:
		\[
		F\land V = F
		\]
		por lo tanto:
		\[
		[F\lor V]\to V
		\]
		
		Entonces:
		\[
		V\to V = V
		\]
		
		\Resultado{
			La proposición compuesta es \textbf{verdadera}.
		}
	\end{frame}
	
	% =========================================================
	\section{Problema 2}
	% =========================================================
	
	% ---------------------------------------------------------
	\begin{frame}{Problema 2: Enunciado}
		Considere la proposición:
		
		\begin{quote}
			``Si una persona no sigue las estrategias terapéuticas indicadas, entonces presenta dificultades emocionales o no logra mejorar su bienestar psicológico.''
		\end{quote}
		
		\begin{enumerate}[label=\alph*)]
			\item Identifique las proposiciones simples y exprese el enunciado en lenguaje simbólico.
			\item Construya una tabla de verdad e indique si corresponde a una tautología, contradicción o contingencia.
		\end{enumerate}
	\end{frame}
	
	% ---------------------------------------------------------
	\begin{frame}{Problema 2(a): Simbolización}
		Definimos:
		
		\begin{align*}
			p &: \text{La persona sigue las estrategias terapéuticas indicadas.} \\
			q &: \text{La persona presenta dificultades emocionales.} \\
			r &: \text{La persona logra mejorar su bienestar psicológico.}
		\end{align*}
		
		Entonces:
		\begin{itemize}
			\item ``no sigue las estrategias'' corresponde a \(\sim p\),
			\item ``dificultades emocionales o no mejora'' corresponde a \(q\lor \sim r\).
		\end{itemize}
		
		\Resultado{
			La proposición en lenguaje simbólico es:
			\[
			\sim p \to (q\lor \sim r)
			\]
		}
	\end{frame}
	
	% ---------------------------------------------------------
	\begin{frame}{Problema 2(b): Tabla de verdad}
		\small
		\[
		\begin{array}{c|c|c|c|c|c}
			p & q & r & \sim p & q\lor \sim r & \sim p \to (q\lor \sim r)\\
			\hline
			V & V & V & F & V & V\\
			V & V & F & F & V & V\\
			V & F & V & F & F & V\\
			V & F & F & F & V & V\\
			F & V & V & V & V & V\\
			F & V & F & V & V & V\\
			F & F & V & V & F & F\\
			F & F & F & V & V & V
		\end{array}
		\]
		\normalsize
		
		\IdeaClave{
			En la última columna aparecen valores verdaderos y falsos. Por eso, la proposición no es tautología ni contradicción.
		}
	\end{frame}
	
	% ---------------------------------------------------------
	\begin{frame}{Problema 2(b): Conclusión}
		\Guion{
			``Al revisar la tabla de verdad, observo que la proposición compuesta toma valor verdadero en algunas filas y valor falso en otra. Por lo tanto, se clasifica como una contingencia.''}
		
		\Resultado{
			\[
			\sim p \to (q\lor \sim r)
			\]
			es una \textbf{contingencia}.
		}
	\end{frame}
	
	% =========================================================
	\section{Problema 3}
	% =========================================================
	
	% ---------------------------------------------------------
	\begin{frame}{Problema 3: Datos del problema}
		Intervenciones terapéuticas:
		\[
		R=\text{Regulación emocional},\qquad
		C=\text{Entrenamiento cognitivo},\qquad
		H=\text{Habilidades sociales}
		\]
		
		A partir de la tabla se obtiene:
		
		\[
		R=\{P1,P3,P5,P7,P9,P11,P13,P15,P17,P19\}
		\]
		
		\[
		C=\{P2,P3,P6,P9,P10,P13,P14,P17,P18\}
		\]
		
		\[
		H=\{P1,P4,P6,P7,P8,P12,P13,P15,P16,P18,P20\}
		\]
	\end{frame}
	
	% ---------------------------------------------------------
	\begin{frame}{Problema 3(a): Organización para el diagrama de Venn}
		\small
		\begin{multicols}{2}
			\begin{block}{Solo \(R\)}
				\[
				\{P5,P11,P19\}
				\]
			\end{block}
			
			\begin{block}{Solo \(C\)}
				\[
				\{P2,P10,P14\}
				\]
			\end{block}
			
			\begin{block}{Solo \(H\)}
				\[
				\{P4,P8,P12,P16,P20\}
				\]
			\end{block}
			
			\begin{block}{\(R\cap C\) solo}
				\[
				\{P3,P9,P17\}
				\]
			\end{block}
			
			\begin{block}{\(R\cap H\) solo}
				\[
				\{P1,P7,P15\}
				\]
			\end{block}
			
			\begin{block}{\(C\cap H\) solo}
				\[
				\{P6,P18\}
				\]
			\end{block}
			
			\begin{alertblock}{Intersección triple}
				\[
				R\cap C\cap H=\{P13\}
				\]
			\end{alertblock}
		\end{multicols}
		\normalsize
	\end{frame}
	
	% ---------------------------------------------------------
	\begin{frame}{Problema 3(a): Guion para explicar el diagrama}
		\Guion{
			``Para construir el diagrama de Venn, distribuyo a las personas según participen en una sola intervención, en exactamente dos o en las tres.\\[0.2cm]
			En solo Regulación Emocional ubico \(P5, P11, P19\).\\
			En solo Entrenamiento Cognitivo ubico \(P2, P10, P14\).\\
			En solo Habilidades Sociales ubico \(P4, P8, P12, P16, P20\).\\
			Luego coloco en las intersecciones dobles a \(P3, P9, P17\), \(P1, P7, P15\) y \(P6, P18\), y finalmente en la intersección triple ubico a \(P13\).''}
	\end{frame}
	
	% ---------------------------------------------------------
	\begin{frame}{Problema 3(b): Exactamente dos intervenciones}
		Se pide:
		\[
		A=\{x\in U \mid x \text{ participa exactamente en dos intervenciones}\}
		\]
		
		Por lo tanto, tomamos solo las intersecciones dobles sin incluir la triple:
		
		\[
		R\cap C \text{ solo }=\{P3,P9,P17\}
		\]
		\[
		R\cap H \text{ solo }=\{P1,P7,P15\}
		\]
		\[
		C\cap H \text{ solo }=\{P6,P18\}
		\]
		
		Entonces:
		\[
		A=\{P1,P3,P6,P7,P9,P15,P17,P18\}
		\]
		
		\Resultado{
			\[
			|A|=8
			\]
		}
	\end{frame}
	
	% ---------------------------------------------------------
	\begin{frame}{Problema 3(c): Conjunto por extensión}
		Se pide el conjunto de personas que:
		
		\begin{center}
			\emph{participan solo en Habilidades Sociales, o participan simultáneamente en Regulación Emocional y Entrenamiento Cognitivo.}
		\end{center}
		
		Tenemos:
		\[
		\text{Solo }H=\{P4,P8,P12,P16,P20\}
		\]
		y
		\[
		R\cap C=\{P3,P9,P13,P17\}
		\]
		
		Por lo tanto:
		\[
		\{P4,P8,P12,P16,P20\}\cup\{P3,P9,P13,P17\}
		\]
		
		\Resultado{
			\[
			\{P3,P4,P8,P9,P12,P13,P16,P17,P20\}
			\]
		}
	\end{frame}
	
	% =========================================================
	\section{Problema 4}
	% =========================================================
	
	% ---------------------------------------------------------
	\begin{frame}{Problema 4: Datos del problema}
		Sea el universo:
		\[
		U=\{1,2,3,\dots,20\}
		\]
		
		Conjuntos dados:
		\[
		A=\{2,4,6,8,10,12,14,16,18,20\}
		\]
		\[
		S=\{5,7,9,11,13\}
		\]
		\[
		E=\{3,6,9,12,15,18\}
		\]
		
		Se pide resolver operaciones entre conjuntos.
	\end{frame}
	
	% ---------------------------------------------------------
	\begin{frame}{Problema 4(a): Desarrollo}
		Calcular:
		\[
		(A-E)\cup (S\cap E)
		\]
		
		Primero:
		\[
		A-E=\{2,4,8,10,14,16,20\}
		\]
		
		porque de \(A\) se eliminan los elementos comunes con \(E\): \(6,12,18\).
		
		Luego:
		\[
		S\cap E=\{9\}
		\]
		
		Finalmente:
		\[
		(A-E)\cup (S\cap E)=\{2,4,8,9,10,14,16,20\}
		\]
		
		\Resultado{
			\[
			(A-E)\cup (S\cap E)=\{2,4,8,9,10,14,16,20\}
			\]
		}
	\end{frame}
	
	% ---------------------------------------------------------
	\begin{frame}{Problema 4(a): Guion sugerido}
		\Guion{
			``Primero calculo \(A-E\), es decir, los elementos de \(A\) que no pertenecen a \(E\).\\
			Como \(A=\{2,4,6,8,10,12,14,16,18,20\}\) y \(E=\{3,6,9,12,15,18\}\), elimino de \(A\) los elementos \(6\), \(12\) y \(18\).\\
			Así obtengo \(A-E=\{2,4,8,10,14,16,20\}\).\\[0.2cm]
			Luego calculo \(S\cap E\), o sea, los elementos comunes entre \(S\) y \(E\). El único común es \(9\).\\
			Por eso \(S\cap E=\{9\}\).\\[0.2cm]
			Finalmente uno ambos resultados y obtengo \(\{2,4,8,9,10,14,16,20\}\).''}
	\end{frame}
	
	% ---------------------------------------------------------
	\begin{frame}{Problema 4(b): Desarrollo}
		Interpretamos la expresión como:
		\[
		\big((A-S)^c\cap A^c\big)-E
		\]
		
		Primero, como \(A\) y \(S\) no tienen elementos en común:
		\[
		A-S=A
		\]
		
		Entonces:
		\[
		(A-S)^c=A^c
		\]
		
		Ahora calculamos el complemento de \(A\) respecto de \(U\):
		\[
		A^c=\{1,3,5,7,9,11,13,15,17,19\}
		\]
		
		Luego:
		\[
		(A-S)^c\cap A^c=A^c\cap A^c=A^c
		\]
		
		Finalmente:
		\[
		A^c-E=\{1,3,5,7,9,11,13,15,17,19\}-\{3,6,9,12,15,18\}
		\]
		\[
		=\{1,5,7,11,13,17,19\}
		\]
		
		\Resultado{
			\[
			\big((A-S)^c\cap A^c\big)-E=\{1,5,7,11,13,17,19\}
			\]
			y su cardinalidad es
			\[
			7.
			\]
		}
	\end{frame}
	
	% ---------------------------------------------------------
	\begin{frame}{Problema 4(b): Guion sugerido}
		\Guion{
			``Como los conjuntos \(A\) y \(S\) no comparten elementos, entonces \(A-S=A\).\\
			Por ello, \((A-S)^c=A^c\).\\
			El complemento de \(A\) en el universo es \(\{1,3,5,7,9,11,13,15,17,19\}\).\\
			Al intersectarlo nuevamente con \(A^c\), el resultado sigue siendo \(A^c\).\\
			Después resto el conjunto \(E\), es decir, elimino \(3\), \(9\) y \(15\).\\
			Así obtengo el conjunto final \(\{1,5,7,11,13,17,19\}\), cuya cardinalidad es \(7\).''}
	\end{frame}
	
	% =========================================================
	\section{Cierre}
	% =========================================================
	
	% ---------------------------------------------------------
	\begin{frame}{Resumen final de respuestas}
		\small
		\begin{block}{Problema 1}
			\begin{itemize}
				\item[a)] Traducción correcta de \((p\to q)\land(\sim r\to s)\).
				\item[b)] \(\;[(p\land \sim q)\lor r]\to s\) es \textbf{verdadera}.
			\end{itemize}
		\end{block}
		
		\begin{block}{Problema 2}
			\begin{itemize}
				\item[a)] \(\sim p\to(q\lor \sim r)\)
				\item[b)] Es una \textbf{contingencia}.
			\end{itemize}
		\end{block}
		
		\begin{block}{Problema 3}
			\[
			A=\{P1,P3,P6,P7,P9,P15,P17,P18\},\qquad |A|=8
			\]
			\[
			\{P3,P4,P8,P9,P12,P13,P16,P17,P20\}
			\]
		\end{block}
		
		\begin{block}{Problema 4}
			\[
			(A-E)\cup(S\cap E)=\{2,4,8,9,10,14,16,20\}
			\]
			\[
			\big((A-S)^c\cap A^c\big)-E=\{1,5,7,11,13,17,19\}
			\]
			Cardinalidad: \(7\)
		\end{block}
		\normalsize
	\end{frame}
	
	% ---------------------------------------------------------
	\begin{frame}{Cierre del video}
		\Guion{
			``En esta evaluación trabajamos lógica proposicional, traducción al lenguaje usual, lenguaje simbólico, tablas de verdad, diagramas de Venn, operaciones entre conjuntos y cardinalidad.\\[0.2cm]
			En cada problema se mostró el desarrollo paso a paso y la justificación matemática correspondiente. Muchas gracias.''}
		
		\vfill
		\begin{center}
			{\Large \textbf{Fin del video}}
		\end{center}
	\end{frame}
	
	% =========================================================
\end{document}